\section{Conclusion}

This work presents GIMAN, a comprehensive graph-based deep learning framework for Parkinson's disease precision medicine that uniquely combines predictive accuracy with interpretability. Through systematic integration of clinical, imaging, genetic, and biochemical data, GIMAN addresses three critical clinical needs: identifying distinct progression subtypes to enable trial enrichment (38\% sample size reduction), stratifying prodromal individuals by conversion risk (C-index 0.79), and providing interpretable explanations validated across six independent methods (91\% cross-method consensus).

Our end-to-end pipeline demonstrates that rigorous data preprocessing, graph-based multimodal learning, and comprehensive explainability validation can collectively produce clinical AI systems worthy of physician trust. The validation that model predictions are grounded in established PD neurobiology—dopaminergic neurodegeneration, $\alpha$-synuclein pathology, genetic risk factors—confirms that GIMAN learned clinically meaningful patterns rather than dataset artifacts. Counterfactual analysis translates predictions into actionable insights by quantifying feature changes required to alter outcomes, providing clear targets for therapeutic intervention.

The broader implications extend beyond Parkinson's disease. First, patient similarity graphs that leverage relationships between individuals offer advantages over traditional approaches treating patients as independent samples, enabling a form of analogical reasoning similar to expert clinical decision-making. Second, multi-method explainability validation substantially increases model trust compared to single-method explanations, establishing a new standard for clinical AI validation. Third, comprehensive frameworks integrating preprocessing through explanation are more suitable for clinical deployment than isolated algorithmic innovations.

As the field advances toward precision medicine, tools that balance predictive power with interpretability become essential. GIMAN demonstrates that these goals are not mutually exclusive—graph-based attention mechanisms provide both strong performance and inherent explainability. With prospective clinical validation and implementation in real-world care settings, GIMAN could transform PD management by enabling truly personalized treatment strategies guided by predicted progression trajectories and explained through transparent, biologically grounded reasoning.

The framework presented here establishes a blueprint for developing trustworthy clinical AI: start with rigorous data curation, leverage structured patient relationships through graphs, apply state-of-the-art deep learning, and validate interpretability through convergent evidence from complementary explanation methods. As disease-modifying therapies emerge for neurodegenerative diseases, such precision medicine tools will be crucial for identifying optimal treatment candidates, monitoring therapeutic responses, and accelerating clinical trial success. GIMAN provides a foundation for this future of data-driven, interpretable, and ultimately more effective neurological care.
